\documentclass{article}
\usepackage{geometry}
\geometry{a4paper,scale=0.8}
\usepackage[utf8]{inputenc}

\usepackage{amsmath, amssymb}
\usepackage{amsthm}
\usepackage{enumitem}

\newtheoremstyle{breakstyle}
    {5pt}
    {10pt}
    {\normalfont}
    {0pt}
    {\bfseries}
    {\newline\indent}
    {0pt}
    {\thmname{#1}\thmnumber{ #2}\thmnote{ \textbf{[#3]}}}
\theoremstyle{breakstyle}
\newtheorem{exercise}{Exercise}[section]
\newtheorem{remark}{Remark}[section]

\title{Chapter 2: Algebraic Numbers}
\author{Flandre}
\date{\today}

\begin{document}
\maketitle

\setcounter{section}{2}

\begin{exercise}
    \leavevmode \vspace{-\baselineskip}
    \begin{enumerate}[label=\alph*)]
        \item Algebraic, not integer.
        \item Algebraic integer.
        \item Algebraic integer.
        \item Algebraic integer.
        \item Algebraic, not integer.
        \item Algebraic integer.
    \end{enumerate}
\end{exercise}

\begin{exercise}
    See \emph{Example 2.3}, we can easily get $\mathbb{Q}(\sqrt{3}, \sqrt[3]{5}) = \mathbb{Q}(\sqrt{3} + \sqrt[3]{5})$.
\end{exercise}

\begin{exercise}
    Notice the minimal polynomial of $\sqrt[3]{7}$ over $\mathbb{Q}$ is $x^3 - 7$, which have roots
    \[ 
        \sqrt[3]{7}, \sqrt[3]{7}\omega, \sqrt[3]{7}\omega^2, \text{ where } \omega \text{ is the } 3\text{rd root of unity}. 
    \]
    Therefore, the 3 monomorphisms will be
    \[ 
        \begin{cases}
            \sigma_1 : \sigma_1(a) = a, \sigma_1(\sqrt[3]{7}) = \sqrt[3]{7} \\
            \sigma_2 : \sigma_2(a) = a, \sigma_2(\sqrt[3]{7}) = \sqrt[3]{7}\omega \\
            \sigma_3 : \sigma_3(a) = a, \sigma_3(\sqrt[3]{7}) = \sqrt[3]{7}\omega^2
        \end{cases}, \text{ where } a \text{ is an arbitary number in } \mathbb{Q}.
    \]
\end{exercise}

\begin{exercise}
    Let monomorphisms of $\mathbb{Q}(\sqrt{3}, \sqrt{5}) = \mathbb{Q}(\sqrt{3} + \sqrt{5}) \longrightarrow \mathbb{C}$ be $\sigma_k, k \in \{1, 2, 3, 4\}$ which holds its value in $\mathbb{Q}$ and
    \[ 
        \begin{cases}
            \sigma_1 : \sigma_1(\sqrt{3} + \sqrt{5}) = \sqrt{3} + \sqrt{5} \\
            \sigma_2 : \sigma_2(\sqrt{3} + 2\sqrt{5}) = -\sqrt{3} - \sqrt{5} \\
            \sigma_3 : \sigma_3(\sqrt{3} + 2\sqrt{5}) = \sqrt{3} - \sqrt{5} \\
            \sigma_4 : \sigma_4(\sqrt{3} + 2\sqrt{5}) = -\sqrt{3} + 2\sqrt{5}
        \end{cases} 
    \]
    Hence, the discriminant of $\mathbb{Q}(\sqrt{3} + \sqrt{5})$ is
    \[ 
        \det \begin{bmatrix}
            1 & \sqrt{3} + \sqrt{5}  & 8 + 2\sqrt{15} & 18\sqrt{3} + 14\sqrt{5} \\
            1 & -\sqrt{3} - \sqrt{5} & 8 + 2\sqrt{15} & -18\sqrt{3} - 14\sqrt{5} \\
            1 & \sqrt{3} - \sqrt{5}  & 8 - 2\sqrt{15} & 18\sqrt{3} - 14\sqrt{5} \\
            1 & -\sqrt{3} + \sqrt{5} & 8 - 2\sqrt{15} & -18\sqrt{3} + 14\sqrt{5}
        \end{bmatrix} = 1920 = 2^7 \times 3 \times 5 
    \]
    Let $\sqrt{3} + \sqrt{5} = \theta$. Therefore, we have to examine whether there exist $\alpha$ such that $\frac{1}{2}(\lambda_1 + \lambda_2 \theta + \lambda_3 \theta^2 + \lambda_4 \theta^3)$ is an algebraic integer where $\lambda_k \in \{0, 1, 2, 3\}$. We first check the norm of it.
\end{exercise}

\begin{exercise}
    \leavevmode \vspace{-\baselineskip}
    \begin{enumerate}[label=\roman*.]
        \item Since the minimal polynomial is $x^4 - 2$, whose roots are $\sqrt[4]{2}, \omega\sqrt[4]{2}, \omega^2\sqrt[4]{2}, \omega^3\sqrt[4]{2}$ where $\omega$ is the 4th root of the unity. The 4 monomorphisms $\sigma_k$ are those holding the value on $\mathbb{Q}$ and
        \[ 
            \begin{cases}
                \sigma_1 : \sigma_1(\sqrt[4]{2}) = \sqrt[4]{2} \\
                \sigma_2 : \sigma_2(\sqrt[4]{2}) = \omega\sqrt[4]{2} \\
                \sigma_3 : \sigma_3(\sqrt[4]{2}) = \omega^2\sqrt[4]{2} \\
                \sigma_4 : \sigma_4(\sqrt[4]{2}) = \omega^3\sqrt[4]{2}
            \end{cases} 
        \]
        \item As we know from \emph{i)}, the minimal polynomial is $x^4 - 2$.
        \item
        \begin{enumerate}[label=\alph*)]
            \item $\prod_{k=1}^4 (x - \sigma_k(\sqrt[4]{2})) = x^4 - 2$
            \item $\prod_{k=1}^4 (x - \sigma_k(\sqrt{2})) = (x^2 - 2)^2$
            \item $\prod_{k=1}^4 (x - \sigma_k(2)) = (x - 2)^4$
            \item $\prod_{k=1}^4 (x - \sigma_k(\sqrt{2} + 1)) = (x^2 - 2x - 1)^2$
        \end{enumerate}
    \end{enumerate}
\end{exercise}

\begin{exercise}
    When $p = 3$, we follow from the book by checking the trace first:
    \[ 
        \mathrm{T}(\alpha) = \lambda_1 \in \mathbb{Z}, \text{ which is useless} 
    \]
    We then compute the norm:
    \begin{align*}
        \mathrm{N}(a + b\theta + c\theta^2) &= (a + b\theta + c\theta^2)(a + b\omega\theta + c\omega^2\theta^2)(a + b\omega^2\theta + c\omega\theta^2) \\
        &= a^3 + 5b^3 + 25c^3 - 15abc \\
        &\equiv 0 \pmod{27}
    \end{align*}
    Hence,
\end{exercise}

\begin{exercise}
\end{exercise}

\begin{exercise}
    \leavevmode \vspace{-\baselineskip}
    \begin{enumerate}[label=\alph*)]
        \item $\mathbb{Q}(\sqrt{2}, \sqrt{3}) = \mathbb{Q}(\sqrt{2} + \sqrt{3})$
    \end{enumerate}
\end{exercise}

\begin{exercise}
\end{exercise}

\begin{exercise}
    Suppose it is not an integral basis, and denote the additive group $\alpha_1, \ldots, \alpha_n$ generated by $G$. Therefore $G$ must be a subgroup of $\mathfrak{O}$, and its discriminant cannot be equal to $d$.
\end{exercise}

\begin{exercise}
    The problem is trivial since any monomorphisms in the problem holds the value in $\mathbb{Q}$.
\end{exercise}

\begin{exercise}
    Let
    \[ 
        \begin{cases}
            K_1 = \mathbb{Q}(x^6 - 1) \\
            K_2 = \mathbb{Q}(x^3 - 1)
        \end{cases} 
    \]
    In this case, from \emph{Exercise 11} we have
    \[ 
        \begin{cases}
            \mathrm{N}_{K_1}(1) = 6 \\
            \mathrm{N}_{K_2}(1) = 3
        \end{cases} 
    \]
\end{exercise}

\begin{exercise}
    \leavevmode \vspace{-\baselineskip}
    \begin{enumerate}[label=\roman*.]
        \item The $n$ monomorphisms have the same forms and proofs as those in \emph{Theorem 2.4}.
        
        \item $\mathrm{N}_{K/L}(\alpha_1 \alpha_2) = \prod_{k=1}^n \sigma_k(\alpha_1 \alpha_2) = \prod_{k=1}^n \sigma_k(\alpha_1) \prod_{k=1}^n \sigma_k(\alpha_2) = \mathrm{N}_{K/L}(\alpha_1) \mathrm{N}_{K/L}(\alpha_2)$
        
        $\mathrm{T}_{K/L}(\alpha_1 + \alpha_2) = \sum_{k=1}^n \sigma_k(\alpha_1 + \alpha_2) = \sum_{k=1}^n \sigma_k(\alpha_1) + \sum_{k=1}^n \sigma_k(\alpha_2) = \mathrm{T}_{K/L}(\alpha_1) + \mathrm{T}_{K/L}(\alpha_2)$
        \qed
        
        \item
        
        \item \emph{(I suppose the author means} $\mathrm{N}_{K/L}(\alpha)$ \emph{since} $\sqrt{\alpha}$ \emph{is not in} $K$\emph{)}
        
        Field $K$ has minimal polynomial $x^2 - \sqrt{3}$, which has roots $\sqrt[4]{3}, -\sqrt[4]{3}$, over $\mathbb{Q}$.
        \begin{enumerate}[label=\roman*.]
            \item When $\alpha = \sqrt[4]{3}$, $\mathrm{N}_{K/L}(\alpha) = \sqrt[4]{3}(-\sqrt[4]{3}) = -\sqrt{3}$, $\mathrm{T}_{K/L}(\alpha) = \sqrt[4]{3} - \sqrt[4]{3} = 0$.
            
            \item When $\alpha = \sqrt[4]{3} + \sqrt{3}$, $\mathrm{N}_{K/L}(\alpha) = (\sqrt[4]{3} + \sqrt{3})(-\sqrt[4]{3} + \sqrt{3}) = 3 - \sqrt{3}$, $\mathrm{T}_{K/L}(\alpha) = (\sqrt[4]{3} + \sqrt{3}) + (-\sqrt[4]{3} + \sqrt{3}) = 2\sqrt{3}$.
        \end{enumerate}
    \end{enumerate}
\end{exercise}

\begin{exercise}
    \leavevmode \vspace{-\baselineskip}
    \begin{enumerate}[label=\roman*.]
        \item $\mathrm{N}_{K/L}(\sqrt{3}) = \sqrt{3} \sqrt{3} = 3$, $\mathrm{N}_{K/\mathbb{Q}}(\alpha) = \sqrt{3}(-\sqrt{3})\sqrt{3}(-\sqrt{3}) = 9$.
        
        \item $\mathrm{T}_{K/L}(\sqrt{3}) = \sqrt{3} + \sqrt{3} = 2\sqrt{3}$, $\mathrm{T}_{K/\mathbb{Q}}(\alpha) = \sqrt{3} + (-\sqrt{3}) + \sqrt{3} + (-\sqrt{3}) = 0$.
    \end{enumerate}
\end{exercise}

\end{document}
